% ----------------------------------------------------------------------
% Preâmbulo — Capacitação de MySQL
% ----------------------------------------------------------------------
% ATENÇÃO: compilar com XeLaTeX. A fonte Lexend é carregada de arquivos
% locais em assets/fonts/, o que exige fontspec (XeLaTeX ou LuaLaTeX).
% ----------------------------------------------------------------------
\usepackage[brazil]{babel}

% --- Fonte da marca (Lexend, local) -----------------------------------
\usepackage{fontspec}
\newcommand{\fontpath}{./assets/fonts/}
\setmainfont[
    Path          = \fontpath,
    UprightFont   = *-Regular,
    BoldFont      = *-Bold,
    AutoFakeSlant = 0.2,          % Lexend não possui itálico real
]{Lexend}

% --- Estrutura e tipografia -------------------------------------------
\usepackage{microtype}
\usepackage{indentfirst}
\usepackage{graphicx}
\usepackage{xcolor}
\usepackage{float}
\usepackage{fancyhdr}

% --- Tabelas -----------------------------------------------------------
\usepackage{longtable}
\usepackage{tabularx}
\usepackage{booktabs}
\usepackage{array}
\usepackage{multirow}

% --- Código ------------------------------------------------------------
\usepackage{listings}

% --- Símbolos usados nos avisos ---------------------------------------
\usepackage{amssymb}
\usepackage{pifont}

% --- Caixas de destaque ------------------------------------------------
\usepackage[most]{tcolorbox}

% --- Links e citações --------------------------------------------------
\usepackage{url}
\usepackage{hyperref}
\usepackage{abntex2cite}

% --- Justificação -----------------------------------------------------
% Nomes de comandos em \texttt{} são longos e não hifenizam, o que empurra
% linhas para fora da margem. O emergencystretch autoriza o TeX a esticar
% mais os espaços nesses parágrafos em vez de estourar a caixa.
\setlength{\emergencystretch}{3em}
